\section{Core Laws}

\begin{law}[Removal idempotence]
\label{law:remove-idempotent}
\[
  \Remove{d}(\Remove{d}(H))=\Remove{d}(H).
\]
After every line containing $d$ has been removed, another removal has no
further effect.
\end{law}

\begin{law}[Commutativity of distinct removals]
\label{law:remove-commute}
For distinct DNS names $d$ and $e$,
\[
  \Remove{d}(\Remove{e}(H))=\Remove{e}(\Remove{d}(H)).
\]
The law remains true when one line contains both names; that line is absent
from either final result.
\end{law}

\begin{law}[Preservation of unrelated lines]
\label{law:preservation}
If $\neg\Contains(d,l)$, then
\[
  \Count(l,\Remove{d}(H))=\Count(l,H).
\]
The law covers blank lines, comments, and tokens that merely contain similar
text.
\end{law}

\begin{law}[Insertion idempotence]
\label{law:insert-idempotent}
\[
  \Insert{d}(\Insert{d}(A))=\Insert{d}(A).
\]
\end{law}

\begin{law}[Setwise commutativity of insertion]
\label{law:insert-commute}
For distinct $d$ and $e$,
\[
  \SetOf{\Insert{d}(\Insert{e}(A))}
  =\SetOf{\Insert{e}(\Insert{d}(A))}.
\]
Sequence equality does not generally hold.  If both names are initially
absent, the two sides append them in opposite orders.
\end{law}

\begin{law}[Exactness]
\label{law:exactness}
For distinct DNS tokens $d$ and $e$,
\[
  e\in\Tokens(l)\land d\notin\Tokens(l)
  \quad\Longrightarrow\quad
  l\text{ occurs unchanged in }\Remove{d}(H).
\]
This is the formal reason that removing \code{ads.example.com} must preserve a
line containing \code{notads.example.com}.
\end{law}

The first four laws admit direct structural proofs from the definitions in the
previous section.  They also translate naturally into property-based tests.
